\documentclass[11pt]{article}
\usepackage[utf8]{inputenc}
\usepackage[margin=1in]{geometry}
\usepackage{graphicx}
\usepackage{float}
\usepackage{xcolor}
\usepackage{booktabs}
\usepackage{caption}
\usepackage{makecell}
\usepackage{enumitem}
\usepackage{hyperref}

\title{\textbf{LOGage 2026: Round 2 — Part 2 Summary}}
\author{Omni-Channel Strategy Analysis}
\date{\today}

\begin{document}
\maketitle
\tableofcontents
\newpage

\section{PART 2: OMNI-CHANNEL FULFILLMENT STRATEGY DESIGN}

\subsection{Question 2.1: Network Model Evaluation \& Design}

\subsubsection{Current Two-RDC Model Assessment}

The brand currently operates two Regional Distribution Centers (RDCs):
\begin{itemize}
  \item \textbf{My Phuoc Warehouse} — Industrial Zone, Binh Duong Province
  \item \textbf{Vinh Loc Warehouse} — Binh Chanh District, Ho Chi Minh City
\end{itemize}

\textbf{Strengths of the current model:}
\begin{itemize}
  \item Good regional separation: My Phuoc serves northern (B2B factory-sourced) flows; Vinh Loc is positioned for inner-city B2C proximity.
  \item Both RDCs handle mixed-channel flows, demonstrating operational flexibility.
  \item Physical locations bracket the Eastern\textendash{}Western commercial corridor of HCMC, covering most \DJ{}\^ong Nam B\^o volume.
\end{itemize}

\textbf{Important context: Non-overlapping operating timelines}
\begin{itemize}
  \item \textbf{My Phuoc} (Primary RDC: Jun to Nov 2025 (6 months)): The \textit{primary} warehouse, operational across the full 6-month analysis window. All June--November 2025 throughput flowed through My Phuoc.
  \item \textbf{Vinh Loc} (New RDC: Dec 2025 only (1 month)): A \textit{new} warehouse that came online in December 2025 only. Raw volume totals are therefore \textbf{not directly comparable} between the two facilities.
\end{itemize}

\textbf{Data-driven channel flow analysis} (per-day rates, Figure \ref{fig:q21-channel-flow}):
\begin{itemize}
  \item \textbf{My Phuoc} (3,492 B2B + 413 B2C orders over 6 months): B2B is a \textit{Mixed flow} accounting for \textbf{99.2\%} of CBM. B2B avg \textbf{25.5 orders/active day}, volatility index \textbf{B2B = 1.79}.
  \item \textbf{Vinh Loc} (1,250 B2B + 47 B2C orders in 1 month): B2B is a \textit{Large-volume flow} accounting for \textbf{99.7\%} of CBM. B2B avg \textbf{48.7 orders/active day}, volatility index \textbf{B2B = 1.94}.
  \item On a per-day basis, Vinh Loc processed B2B orders at a comparable daily rate to My Phuoc despite being open for only one month, indicating strong December surge demand (year-end push).
\end{itemize}

\begin{figure}[H]
\centering
\includegraphics[width=\linewidth]{../charts/q21_channel_flow_profile.png}
\caption{Q2.1 Warehouse $\times$ Channel Flow Profile (per-day rates normalized to each RDC operating window)}
\label{fig:q21-channel-flow}
\end{figure}

\textbf{Limitations for simultaneous B2B + B2C:}
\begin{itemize}
  \item Lead times from both RDCs are too long for the 2\textendash{}4 hour B2C SLA required for e-commerce \textendash{} even Vinh Loc is \textgreater{}10 km from central districts.
  \item Shared inventory creates channel conflict: B2B bulk orders may deplete stock that B2C needs for same-day fulfillment.
  \item Pick-and-pack processes for large pallet B2B orders disrupt small-unit B2C wave picking.
\end{itemize}

\subsubsection{B2C E-Commerce SLA Assessment}

To meet a 2--4 hour delivery SLA across Ho Chi Minh City, we applied a 25 km distance threshold from the nearest RDC as the feasibility cut-off (assuming average urban speed of $\sim$25 km/h, 25 km $\approx$ 60 min travel, leaving buffer for pick/pack and handoff).

The analysis of \textbf{21 HCM districts} reveals a clear split:
\begin{itemize}
  \item \textbf{\color{green!60!black}21 districts — Adequate:} These districts can be served within the 2h SLA from the existing RDCs and account for \textbf{100.0\%} of total HCM B2B volume (\textbf{65,527 units}).
  \item \textbf{\color{red}0 districts — Require Dark Stores:} These are \textgreater{}25 km from both RDCs, placing them outside reliable SLA reach. They represent \textbf{0.0\%} of volume (\textbf{0 units}).
\end{itemize}

Figure \ref{fig:q21-network-coverage} shows the district-level breakdown.

\begin{figure}[H]
\centering
\includegraphics[width=\linewidth]{../charts/q21_network_coverage.png}
\caption{Q2.1 HCM District Proximity to Existing RDCs vs B2B Demand Volume}
\label{fig:q21-network-coverage}
\end{figure}

\subsubsection{Dark Store Node Recommendation}

Based on the demand concentration and distance analysis, we recommend placing \textbf{1--2 urban micro-fulfillment (dark store) nodes} in HCMC:
\begin{enumerate}
  \item \textbf{Primary node — Bình Tân / Tân Phú corridor}: These two adjacent districts are among the highest-volume B2B districts (combined \textgreater{}50,000 units) and are already well-served by Vinh Loc (5--7 km). Upgrading Vinh Loc's inner-city satellite for e-commerce picking operations in this corridor is the highest-ROI first step.
  \item \textbf{Secondary node — District 1 / District 3 / Phú Nhuận cluster}: The CBD and premium residential cluster lies 14+ km from Vinh Loc, making 2h SLA difficult. A small dark store (200--300 m²) in this zone, stocked with the top Class A SKUs, would close the gap.
\end{enumerate}

\subsubsection{Order Split Logic Design (HCM)}

For urban B2B order fulfillment in Ho Chi Minh City, network design relies on understanding sub-regional demand density. We analyzed the historical B2B shipments across all HCM districts.

The demand is highly concentrated: the top 5 districts alone command \textbf{56,356 units} — \textbf{86.0\%} of all HCM B2B volume — as shown in Figure \ref{fig:q21-hcm-district-volume}.

\begin{figure}[H]
\centering
\includegraphics[width=0.85\linewidth]{../charts/q21_hcm_district_volume.png}
\caption{Q2.1 Top Ho Chi Minh Districts by B2B Quantity}
\label{fig:q21-hcm-district-volume}
\end{figure}

\textbf{Recommended order split logic:} Given this concentration, the dispatch decision tree should be:
\begin{enumerate}
  \item If delivery district is within 25 km of Vinh Loc $\rightarrow$ dispatch from \textbf{Vinh Loc} (B2C-optimized lane).
  \item If delivery district is within 35 km of My Phuoc but $>$25 km from Vinh Loc $\rightarrow$ dispatch from \textbf{My Phuoc} (B2B bulk lane).
  \item If district is $>$25 km from both RDCs $\rightarrow$ \textbf{route via dark store} (pending network expansion).
\end{enumerate}

\subsection{Question 2.2: Inventory Optimization}

\subsubsection{Lead Time Assumptions \& Validation}

\textbf{Definition:} Replenishment Lead Time is the duration from when the warehouse identifies a replenishment need until the goods are ready to pick again. This is distinct from the customer-facing delivery SLA.

\textbf{Data limitation:} The dataset only contains \texttt{created\_date} (outbound dispatch date). There is no PO creation date or Goods Receipt date, so lead time cannot be measured directly. The values below are \textit{justified assumptions} cross-validated against inter-shipment intervals in the transaction log.

\begin{table}[H]
\centering
\caption{Lead Time Component Breakdown by Channel}
\begin{tabular}{lccl}
\toprule
\textbf{Component} & \textbf{B2B (Factory→RDC)} & \textbf{B2C (RDC→Dark Store)} & \textbf{Rationale} \\
\midrule
T\_review (stock check) & 1.0 day & 0.5 day & B2B daily review; B2C twice-daily auto-trigger \\
T\_PO (create \& approve PO) & 0.5 day & 0 day & B2C uses auto-replenishment, no manual PO \\
T\_supplier (preparation) & 1.5 days & 0 day & Factory batch lead time; B2C pulls from RDC stock \\
T\_transport & 1.0 day & 0.5 day & Intra-ĐNB region; B2C last-mile $<$15 km same-day \\
T\_receiving (QC check-in) & 0.5 day & 0.5 day & Physical count \& inspection \\
T\_putaway (shelf-ready) & 0.5 day & 0.5 day & Ready to pick \\
\midrule
\textbf{Total Lead Time} & \textbf{5 days} & \textbf{2 days} & \\
\bottomrule
\end{tabular}
\end{table}

\textbf{Empirical cross-validation from data:}
\begin{itemize}
  \item Median inter-shipment gap for Class A SKUs = \textbf{2.0 days} $\rightarrow$ matches LT\textsubscript{B2C} = 2 days
  \item Mean inter-shipment gap = \textbf{7.7 days} (long tail from bulk lot orders) $\rightarrow$ consistent with LT\textsubscript{B2B} $\approx$ 5 days
  \item \textit{Conclusion: Assumptions are well-calibrated to observable shipment rhythms.}
\end{itemize}

\subsubsection{Safety Stock Formula \& Results}

Safety stock for each Class A SKU is calculated using the full lead-time uncertainty formula:
\[
SS = Z \cdot \sqrt{LT \cdot \sigma_d^2 + \mu_d^2 \cdot \sigma_{LT}^2}
\]
where $Z = 1.645$ (95\% service level), $\mu_d$ = mean daily demand, $\sigma_d$ = std of daily demand, $LT$ = mean lead time, and $\sigma_{LT}$ = lead time std dev (B2B: 1.0 day, B2C: 0.5 day).

\textbf{Key result:} This section details the required safety stock for Class A SKUs to support omni-channel fulfillment under different lead time configurations.

Under the \textbf{separated channel} setup (B2B: 5-day LT, B2C: 2-day LT), the total required safety stock for Class A SKUs is \textbf{31,852 units}.

We then evaluated the impact of varying the replenishment lead time on total safety stock requirements. As illustrated in Figure \ref{fig:q22-lead-time-sensitivity}, safety stock exhibits a sub-linear (square root) growth relative to lead time.

\begin{figure}[H]
\centering
\includegraphics[width=0.85\linewidth]{../charts/q22_lead_time_sensitivity.png}
\caption{Q2.2 Lead Time Sensitivity for Class A Safety Stock}
\label{fig:q22-lead-time-sensitivity}
\end{figure}

\subsubsection{Inventory Pooling Strategy}

\textbf{Principle — Why pooling reduces safety stock:}
\begin{itemize}
  \item Separated: $SS = Z\sigma(\sqrt{LT_{B2B}} + \sqrt{LT_{B2C}}) = Z\sigma(\sqrt{5}+\sqrt{2}) = Z\sigma \times 3.650$
  \item Pooled: $SS = Z\sigma\sqrt{\frac{LT_{B2B}+LT_{B2C}}{2}} = Z\sigma\sqrt{3.5} = Z\sigma \times 1.871$
  \item \textbf{Reduction factor} = $\frac{\sqrt{3.5}}{\sqrt{5}+\sqrt{2}} = 0.5125$ $\rightarrow$ pooling reduces safety stock by $\approx$48.7\%
\end{itemize}

By centralizing inventory into a shared pool with an effective average lead time of 3.5 days, the required safety stock drops from \textbf{31,852} to \textbf{16,325 units}, saving \textbf{48.7\%} as shown in Figure \ref{fig:q22-inventory-pooling}.

\begin{figure}[H]
\centering
\includegraphics[width=0.75\linewidth]{../charts/q22_inventory_pooling.png}
\caption{Q2.2 Inventory Pooling Impact on Safety Stock}
\label{fig:q22-inventory-pooling}
\end{figure}

\textbf{Dynamic Allocation Rules (to prevent simultaneous stockout):}
\begin{enumerate}
  \item \textbf{R1 — Common Pool:} Merge B2B and B2C inventory per SKU into a single real-time pool tracked in WMS.
  \item \textbf{R2 — Reserved Floor per channel:} Each channel maintains a floor reserve = its own ROP. Neither channel can dip below the other's floor.
  \item \textbf{R3 — Amber Alert:} When total pool $<$ (SS\textsubscript{B2B} + SS\textsubscript{B2C}): system raises alert and prioritizes confirmed B2B purchase orders (contractual commitments).
  \item \textbf{R4 — Red Conflict:} If same SKU is simultaneously short for both channels: B2B receives its full contracted PO quantity (penalty risk); B2C receives remainder + escalated to next pick wave + customer notified.
  \item \textbf{R5 — AZ Spike Freeze:} For AZ-class SKUs showing demand spikes: temporarily freeze B2C allocations, concentrate stock for B2B, trigger emergency replenishment.
  \item \textbf{R6 — End-of-Day Rebalance:} Rebalance channel inventory ratios to target ROP proportions based on rolling 7-day demand forecast.
\end{enumerate}

\end{document}
